% ============================================================
% Bem-vindo ao TexLab!
% Este arquivo e um tutorial basico de LaTeX.
% Cada secao ensina algo novo. Leia os comentarios (linhas
% que comecam com %) para entender o que esta acontecendo.
% ============================================================

% --- 1. PREAMBULO ---
% Tudo que vem ANTES de \begin{document} e o preambulo.
% Aqui voce configura o documento: tipo, pacotes, etc.

\documentclass[12pt, a4paper]{article}
% \documentclass define o tipo de documento.
% Opcoes comuns: article, report, book, beamer (apresentacao).

% Pacotes adicionam funcionalidades ao LaTeX.
\usepackage[utf8]{inputenc}   % Permite acentos e caracteres especiais
\usepackage[T1]{fontenc}      % Codificacao de fontes para portugues
\usepackage{amsmath, amssymb} % Equacoes matematicas avancadas
\usepackage{graphicx}         % Incluir imagens
\usepackage{hyperref}         % Links clicaveis no PDF
\usepackage{booktabs}         % Tabelas mais bonitas
\usepackage{biblatex}         % Gerenciamento de referencias
\addbibresource{references.bib} % Arquivo com as referencias

% Informacoes do documento
\title{Aprenda LaTeX do Zero}
\author{Seu Nome}
\date{\today}


% --- 2. INICIO DO DOCUMENTO ---
\begin{document}

\maketitle  % Gera o titulo com autor e data

% --- 3. RESUMO ---
\begin{abstract}
    Este documento e um tutorial basico de LaTeX.
    Nele voce vai aprender a escrever texto, criar equacoes,
    montar tabelas e gerar referencias bibliograficas.
\end{abstract}


% --- 4. SECOES ---
% Secoes organizam o documento. Use \section, \subsection, etc.

\section{Texto Basico}

No LaTeX, voce escreve texto normalmente no editor.
Para formatar, use comandos entre chaves:

\begin{itemize}
    \item \textbf{Negrito} com \verb|\textbf{...}|
    \item \textit{Italico} com \verb|\textit{...}|
    \item \underline{Sublinhado} com \verb|\underline{...}|
    \item \texttt{Monoespacado} com \verb|\texttt{...}|
\end{itemize}

Para criar um paragrafo novo, basta deixar uma linha em branco
no codigo fonte. O LaTeX cuida do espacamento.


\section{Listas}

Existem dois tipos de listas:

Lista com marcadores (ulletize):
\begin{itemize}
    \item Primeiro item
    \item Segundo item
    \item Terceiro item
\end{itemize}

Lista numerada (enumerate):
\begin{enumerate}
    \item Primeiro passo
    \item Segundo passo
    \item Terceiro passo
\end{enumerate}


\section{Equacoes}

O LaTeX e incrivel para equacoes matematicas.

Uma equacao na linha do texto: $E = mc^2$ (usa \$ para entrar no modo matematico).

Uma equacao centralizada e numerada:
\begin{equation}
    \int_{a}^{b} f(x) \, dx = F(b) - F(a)
    \label{eq:integral}
\end{equation}

A Equacao~\ref{eq:integral} e o Teorema Fundamental do Calculo.

Outras equacoes uteis:
\begin{itemize}
    \item Frações: \verb|\frac{a}{b}| $\rightarrow \frac{a}{b}$
    \item Potencia: \verb|x^2| $\rightarrow x^2$
    \item Raiz: \verb|\sqrt{x}| $\rightarrow \sqrt{x}$
    \item Soma: \verb|\sum_{i=1}^{n} i| $\rightarrow \sum_{i=1}^{n} i$
\end{itemize}


\section{Tabelas}

Tabelas usam o ambiente \texttt{tabular}. Exemplo:

\begin{table}[h]
    \centering
    \begin{tabular}{l c r}
        \toprule
        \textbf{Comando} & \textbf{Resultado} & \textbf{Uso} \\
        \midrule
        \verb|\textbf{}|  & \textbf{Negrito}  & Emphasi      \\
        \verb|\textit{}|  & \textit{Italico}   & Citacao      \\
        \verb|\frac{}{}|  & $\frac{1}{2}$      & Fracao       \\
        \bottomrule
    \end{tabular}
    \caption{Comandos basicos de formatacao}
    \label{tab:comandos}
\end{table}

A Tabela~\ref{tab:comandos} mostra os comandos mais usados.


\section{Imagens}

Para incluir uma imagem, coloque o arquivo na mesma pasta
do \texttt{.tex} e use:

\begin{verbatim}
\begin{figure}[h]
    \centering
    \includegraphics[width=0.5\textwidth]{imagem.png}
    \caption{Descricao da imagem}
    \label{fig:exemplo}
\end{figure}
\end{verbatim}

(Destaque a linha acima com \verb|\begin{verbatim}| para
mostrar o codigo raw. Em um arquivo real, voce usaria
o comando \texttt{\textbackslash includegraphics} normal.)


\section{Referencias Bibliograficas}

Referencias ficam em um arquivo \texttt{.bib} separado.
Para citar no texto, use:

\begin{itemize}
    \item \verb|\cite{chave}| para citar: \cite{knuth1984}
    \item \verb|\parencite{chave}| para citar com parenteses: \parencite{lamport1994}
\end{itemize}

No final do documento, use \verb|\printbibliography| para
gerar a lista de referencias (ela aparece la embaixo).


\section{Dicas Finais}

\begin{enumerate}
    \item \textbf{Salve e compile.} O PDF so atualiza apos compilar.
    \item \textbf{Leia os erros.} O log mostra onde esta o problema.
    \item \textbf{Use comentarios.} Linhas com \% nao aparecem no PDF.
    \item \textbf{Organize com secoes.} Documentos grandes ficam mais faceis de navegar.
    \item \textbf{Nao tenha medo de errar.} E so corrigir e recompilar!
\end{enumerate}

% Gerar a lista de referencias
\printbibliography

\end{document}
